\documentclass[11pt,a4paper]{article}
\usepackage[utf8]{inputenc}
\usepackage[T1]{fontenc}
\usepackage[margin=2.4cm]{geometry}
\usepackage{amsmath,amssymb}
\usepackage{microtype}
\usepackage{enumitem}
\setlist[itemize]{itemsep=2pt,topsep=3pt,leftmargin=1.4em}
\usepackage[hidelinks]{hyperref}
\setlength{\parskip}{5pt}
\setlength{\parindent}{0pt}
\begin{document}
\section*{Response to the Reviewers}

\textbf{Manuscript} 1568909 --- \emph{A 3D Multiphase Continuum Model for Atoms}
\textbf{Revised title} \emph{A 3D Multiphase Continuum Model for Atoms and the Periodic Table}
\textbf{Journal} International Journal of Quantum Chemistry
\textbf{Decision} Major revision

\medskip\hrule\medskip

We thank both reviewers. On the points that mattered most they were correct, and we have not
argued with either: Reviewer 2's objection to the packing argument and Reviewer 1's "elephant in
the room" both identified places where the framing claimed more than the body text supported.

Two things about the revision should be said before the point-by-point replies, and we keep them
short here because both are presented properly in the manuscript.

\textbf{A new account of the build-up of the periodic table} (§4.3). The table is assembled from two
packings the calculations themselves produce --- the antipodal pair, 2, and the octet as two dual
tetrahedra, 8 = 4+4 --- and from one constraint: a row ending in 4 or 6 would need half an octet or
three quarters of one, which is not a packing available to the model, so the remainder after octets
is forced to 0 or 2. That gives every row length 2, 8, 8, 18, 18, 32, 32 and every closed shell
2, 10, 18, 36, 54, 86, 118, with $2n^2$ nowhere in the construction --- so its agreement with the
hydrogenic count is a result, not an input. Three consequences are new: the octet is the unit every
later row is built from (18 a pair and two octets, 32 four octets); the step before closure is 4+3,
so a halogen is short of a tetrahedral \emph{site} rather than of a count; and
$\sqrt{\mathrm{cap}(n{+}1)/\mathrm{cap}(n)}=(n{+}1)/n$ holds exactly when each capacity serves two
periods, which is a direction on the period doubling and is labelled as no more than that. We claim
a decomposition of the capacities, not a derivation of the row lengths; that each capacity serves two
periods remains the principal open item. The 4+4 octet is Linnett's double quartet [\emph{J. Am. Chem. Soc.} \textbf{83} (1961) 2643], which separates the quartets by spin; this model has none and reaches the
same arrangement from Coulomb repulsion and non-overlap alone.

\textbf{Periods 1--4, with no parameters at all.} The energies run hydrogen to krypton, $Z=1$--$36$,
all-electron with the bare term $-Z/r$: no core radius, no screened kernel, no frozen core, no fitted
length, nothing calibrated against a reference atom, nothing set per element. The frozen core and
pseudo-kernel of the original submission, with the softening radius $r_c$ Reviewer 2 identified, are
gone; resolution recovers the range in their place. The only quantity that varies across the table is
the mesh --- $N=400$ through argon, $N=800$ from potassium, since $r_{1s}\approx4.3/Z$ and no single
mesh serves both ends --- and a mesh is a discretisation, not a parameter: it carries no physics and is
fitted to nothing.

\textbf{A word on length.} The manuscript has grown from 13 pages to 28, and we should say plainly why,
since a revision asked for retractions does not usually come back twice the size. The material we
withdrew is substantial --- both heavy-element tables, all the ionization-energy material, the
frozen core and the pseudo-kernel --- and the growth is entirely in what replaced it: the
pair-and-octet construction, the Aufbau section with its measurements, three figures of the
build-up in three dimensions, the fourth-row table, electronegativity as the valence observable,
and the analysis of what the spherical reduction costs. None of it was in the original, and most of
it exists because the reviewers' objections pointed there. \textbf{If the editor considers the paper too
long, we are willing to shorten it} and would propose cutting the discussion of packing versus
orbitals and the more expository passages, leaving the construction, the tables and the figures
intact.

\section*{Reviewer 1}

\textbf{R1.1 --- The ontological status of RealQM. The justification is deferred to a reference that is
itself under review, and the claim that RealQM is a correction to rather than an approximation
of standard quantum mechanics cannot rest on atomic results alone. At minimum the ontological
claims should be phrased as speculation.}

We accept this without reservation and have taken the off-ramp the reviewer offered. Section 1
now closes with a paragraph (\emph{Comparison with textbook quantum mechanics}) stating that the
two are different mathematical objects --- a linear Schrödinger equation for one wave function on
$3N$-dimensional configuration space, versus a nonlinear free-boundary problem for $N$ charge
densities on three-dimensional physical space --- that exclusion is realised differently in each,
and that \textbf{neither is shown here to be an approximation to the other}. It states explicitly
that the wider ontological reading advanced elsewhere should be taken as speculative pending
evidence stronger than atomic energetics can supply, and that nothing in this paper depends on
it. What is compared here is what each produces for the periodic table, and nothing more.

\textbf{R1.2 --- Could equivalence with density functional theory be established?}

Not attempted, and we should say why we do not expect it to be a short step. RealQM is
fundamentally different from DFT in the object it works with. DFT is built on a single common
density rho = sum\_i |psi\_i|\textasciicircum{}2, and Hohenberg-Kohn is a statement about that one function. RealQM
has no such object in that role: its state is N \emph{individual} charge densities together with the
partition of space between them, which carries strictly more information than their sum. Two
different partitions of the same total density have different energies, because the i != j rule
counts pairs over disjoint domains rather than integrating a common rho against itself.

The consequence is visible. DFT's Hartree term integrates rho against itself and so carries a
spurious self-interaction, which the exchange-correlation functional is partly there to remove;
here there is nothing to remove. An equivalence would have to relate objects of different kinds,
and relate an approximated functional to one written down exactly. We name it as open.

\textbf{R1.3 --- Does the method scale to molecules? Is the spherical model useful beyond atoms?}

The spherical reduction is specific to atoms and does not extend; it is available here only
because an atom is spherically symmetric, and the revised introduction now derives the
reduction from that fact rather than presenting it as an implementation choice. The
three-dimensional form of the model carries no such restriction and is the one that extends.

On molecules, what the method supplies first is \textbf{forces}. The free boundaries move with the
charge, so the force on each kernel is available at every step of the same relaxation, with no
separate apparatus and no derivative of a fitted functional; geometry follows from it, and
molecular dynamics follows without further machinery. Energy differences within one composition
are reliable for the same reason, the discretisation error being common to the states compared.
Absolute energetics across compositions is the harder case and we make no claim about it here.
No molecular result is reported in this paper.

\textbf{R1.4 --- Could basis sets reduce the cost, particularly for periodic systems?}

We do not think a basis is needed; the mesh is enough. A basis exists to make the integrals of
a configuration-space method tractable and to compress the representation, and neither problem
arises: there are no two-electron integrals, the interaction coming from a Poisson solve on the
same grid, and the cost is already linear in the number of electrons.

There is also a reason of principle. The unknowns include the interfaces, which are geometric: a
domain boundary is a moving surface, represented directly by a mesh, where a global basis would
have to build a nearly discontinuous partition out of smooth functions. And a basis is an
empirical input whose adequacy must be established per system; introducing one would give back
exactly the parameter-free character R2.1 asks us to secure. Where more accuracy is wanted we
refine the mesh.

\textbf{R1.5 --- Is the imaginary-time propagation connected to path-integral or centroid formalisms?}

No. It is a relaxation scheme used to reach the stationary state and has no such connection.
Only the converged solution enters any result, so the temporal ordering of the relaxation is
immaterial. This is now stated in §Limitations.

\medskip\hrule\medskip

\section*{Reviewer 2}

\textbf{R2.1 --- The abstract claims an in-principle parameter-free calculation, but the results use
per-period $r_c$ calibrations, and separate ones for total energies and ionization energies.}

Correct, and this was the decisive report on the paper. The two statements in the original
manuscript could not both stand. As described above, we did not repair the claim --- we removed
the construction that violated it. There is no $r_c$ anywhere in the revised paper, no
calibration of any kind, and no per-element or per-period empirical input. The inputs are the
nuclear charge and the number of electrons.

What made that possible is worth reporting, because it was not obvious to us. The calibrated
radius existed to let the heavier elements be computed on a fixed mesh; its real function was
to remove the innermost shell, which is the one a uniform grid resolves worst. But that is a
resolution problem and it has a resolution answer. The free-boundary shell radii are outputs of
the model, and we now measure the innermost one to follow $r_{1s}Z \approx 4.3$, constant to
$\pm9\%$ from carbon to argon. Refining the mesh accordingly --- and scaling the step count with
it, since the relaxation is explicit --- the all-electron energies run from argon to zinc within
1--5\% with nothing fitted. The parameter was standing in for a mesh.

(The original abstract also asserted a core radius scaling as $1/Z$ as a \emph{result}. The body
never supported it and it has been withdrawn; what is stated above is a measurement used to
choose a grid, not a claim about the model's output.)

\textbf{R2.2 --- The radial model favours cap-8 packing, yet cap-4 configurations are chosen because
they agree better with reference energies.}

The reviewer's reading was accurate: the packing was selected by fit. The revision replaces that
selection with two things --- a measurement of why the energy cannot decide, and a construction of the
capacities that does not appeal to energies at all.

\textbf{The energy is monotone in the packing, and anti-correlated with the physical configuration.} There
is no interior minimum, so an unconstrained minimiser absorbs the core into the valence. Beryllium's
three candidates order by compactness: a single 4-shell at $-18.893$ Ha (28.8\% too deep), $2{+}2$ at
$-14.893$ (1.54\% too deep), $2{+}1{+}1$ at $-14.603$ (0.44\% short). The lower energy is reached by
over-binding past the reference, so the physical configuration is the least bound and the closest ---
and an energy preference could not have supported it in any case. The $2{+}2$ margin is converged, on
a step ladder flat from 16,000 to 64,000 steps; it is the monotonicity, not a search failure.

\textbf{What decides instead is the relaxation.} The algorithm does not minimise over configurations; it
reaches the equilibrium at the end of its path. A filled shell is more rigid than the charge arriving
at it, and the interface moves only where one domain's amplitude locally exceeds the other's, so the
newcomer is turned aside by a comparison of amplitudes and not by an imposed rule. Figures 2 and 3
show this in three dimensions, all-electron, with no $r_c$: a charge launched as a half-shell wraps a
Li$^+$ core and closes into a shell, declining a merged state that was reachable throughout and lies
1.055 Ha \emph{lower}. A seed ladder rules out the obvious objection that the interface merely stayed
where it was put --- seeds 60\% apart give shell radii 2.139 and 2.153 a.u., 0.7\% apart.

\textbf{The capacities are not fitted either}, and this is the part we would now put forward as the
paper's contribution. They are built from the two filled packings the calculations establish --- the
antipodal pair and the dual-tetrahedral octet, the latter against a single 8-shell that over-binds by
9--22\% --- together with the constraint that a row cannot end in 4 or 6. Every row length and every
closed shell follows, without $2n^2$ entering. §4.3 gives the construction, the 4+3 step before
closure, the statement that only a noble gas can carry a complete outermost octet, and the direction
on the doubling. One classical result is conceded there rather than left for a reader to catch: for
eight points on \emph{one} sphere the Coulomb minimum is the square antiprism, not the cube, by 0.33\%, and
that is a theorem --- though the antiprism is itself a 4+4, and this model does not place eight charges
on one sphere.

What we do \emph{not} claim is that the row lengths are derived. The capacities decompose into pairs and
octets uniquely, and that is established; that the third row has length 8 and the fourth 18 is not.

\textbf{R2.3 --- Equation (1) identifies $-\varepsilon_i$ with the optimised neutral--ion energy
difference. Why should these be equal?}

They are not, and we measured the discrepancy before deciding what to do about it. For
bare-nucleus balanced configurations, Koopmans $-\varepsilon_i$ against $\Delta$SCF against
NIST (eV):

\begin{center}\begin{tabular}{lrrr}\hline
 & Koopmans & $\Delta$SCF & NIST \\ \hline
C & 4.85 & 8.74 & 11.26 \\
O & 6.79 & 12.04 & 13.62 \\
\hline\end{tabular}\end{center}

The frozen estimate is roughly half the relaxed one, so the identification in the original
Eq. (1) was unsupported. Worse for the method, $\Delta$SCF cannot simply replace it: it
requires choosing the cation's packing, and for nitrogen neither choice works --- repacking gives
$+51.4$ eV, decrementing the outer shell gives a state \emph{below} the neutral and a negative
ionization energy. \textbf{$\Delta$SCF therefore inherits the packing-selection problem of R2.2.}

Given this, we concluded that ionization is the wrong observable for a ground-state model:
it is a \emph{process}, in which the remaining electrons repack, and the model has no principled way
to select the ion's packing. \textbf{All ionization-energy material has been removed from the
paper}, and the abstract and Scope now state that ionization and excitation are left aside,
with that reason given.

In its place we report a quantity the model is built to answer: the \textbf{Allen configuration
energy} $\chi = (n_s\varepsilon_s + n_p\varepsilon_p)/(n_s+n_p)$, computed natively from the
shell energies. It is a state property of the neutral atom --- no ion, no anion, no relaxation,
no calibration --- and it is what bond polarity is built from. The \textbf{ordering across period 2 is
exact}. The gradient is too steep by a factor $2.23$ (span $51.1$ against $22.9$ eV), with the
error concentrated at the electronegative end and climbing monotonically from $1.09$ at Be to
$2.06$ at Ne. We present the result where it is strong and state the rest as the known
light-$p$-block limitation, whose structural cause (the outer tetrahedron placed at too large a
radius in a radial model) and cure (the three-dimensional solve) are both identified.

\textbf{R2.4 --- How were the transition-metal points of Figure 1 obtained, and do they establish the
claimed 4s-before-3d removal order?}

They do not. On checking the generating code we found the $s^2$ shell is placed \textbf{outermost by
construction}; the run then tests whether the outermost shell is the $s$ shell, which in a
nested radial model the input very nearly guarantees. The reviewer's suspicion was correct:
4s-before-3d was an input, not a result. Two further defects: those runs used 4000 relaxation
steps, which is below this project's own convergence floor, and their total energies were
floor-carried and not independently meaningful.

\textbf{The claim has been removed from the abstract and from the paper.} The $d$-block is now
outside the stated scope entirely, and §\emph{Boundary of the spherical model} reports the $4s/3d$
interpenetration as the structure the radial reduction cannot represent --- described as a
limitation, with no ordering claimed as a result. §Limitations states that the placement of the
$d$-shell inside the $s$-shell is imposed, not derived.

\textbf{R2.5 --- The Kr→Xe increase contradicts the claimed down-group trend.}

\textbf{R2.6 --- Argon appears as 11.29 eV in Table 3 and 8.0 eV in Table 5.}

\textbf{R2.7 --- $p$-block ionization energies repeat across periods 2--4.}

These three concerned ionization energies of heavy elements and are resolved by the removals
described under R2.1 and R2.3: both source tables are gone from the paper, as is all
ionization-energy material. We record the diagnosis rather than leave it unstated, since each
was a real defect and the reviewer was right about all three:

\begin{itemize}
  \item \textbf{R2.6} was the two calibrations of R2.1 surfacing --- the same atom appearing under the
    total-energy calibration and the ionization-energy calibration without either being labelled.
  \item \textbf{R2.7} was structural, not a transcription error. With the core absorbed into $r_c$ and the
    screened charge identical for every noble-gas valence, each column member presented the \emph{same}
    valence problem, so the entries repeated by construction: the model reproduced the column and
    had no representation of the period. That is a serious limitation of the frozen-core
    construction and is among the reasons we removed it rather than defended it.
  \item \textbf{R2.5} we did not fully diagnose before the construction that produced it was withdrawn.
\end{itemize}

\textbf{R2.8 --- Grid convergence is demonstrated only for the helium total energy.}

The reviewer is right that the ladder covered one atom. We have looked beyond it, and the answer is
not more rungs. Refining the radial mesh converges the calculation to the answer of the \emph{spherically reduced} problem, and that reduction --- not the discretisation --- is what limits the accuracy here. The
revision measures the reduction's cost directly: it has the pair \emph{count} right and the pair
\emph{geometry} wrong, overestimating intra-shell repulsion by a factor of 1.6--2, leaving a residual of
order a per cent in the under-binding direction. Against that, successive $\sqrt2$ refinements of
helium move the total energy by about 0.2\% a step. The model error is an order of magnitude the
larger, so a finer radial grid buys precision in solving the wrong equation. The accuracy-limiting
direction is \emph{angular} resolution, restoring the pair geometry the reduction discards, and the paper
says so instead of offering a mesh ladder as the open question. $N=400$ is also what the
three-dimensional solver reaches, at $400^3$.

\textbf{What the refined runs cost us, reported rather than left to be found.} Argon agrees to $+0.1\%$ at
$N=400$, the closest entry in the table, and it does not survive: on the same domain at $N=724$ and
$N=1024$ it moves to $-511.7$ and $-511.3$ Ha, about 3\% \emph{under}-bound, the two finest rungs agreeing
to 0.08\% and both step-converged. Neon is comparable in size and non-monotone over the same range.
Two changes follow. The paper now states that the table is to be read at the level of a few per cent,
with no single entry claimed to better than a per cent; and we have \textbf{removed the claim that the
model over-binds the heavier third-period atoms with the opposite sign of error to Hartree--Fock},
since the sign of argon's error is a function of the grid. What survives refinement, and what the
reduction predicts, is under-binding --- the direction of all eighteen fourth-row entries.

\section*{Summary of changes}

\textbf{Removed}
\begin{itemize}
  \item All ionization-energy material, including Eq. (1)'s Koopmans identification, and the tables
    underlying R2.5--R2.8.
  \item The frozen core, the pseudo-kernel, $r_c$, and the two-scale core/valence decomposition.
  \item The "4s before 3d" claim (R2.4) and the "in-principle parameter-free" claim as worded (R2.1),
    the latter now restated and true.
  \item The claim that the model over-binds the heavier third-period atoms with the opposite sign of error
    to Hartree--Fock (R2.8): argon's error changes sign between $N=400$ and $N=1024$.
  \item The sentence Reviewer 2 quoted on the correctness of the radial description.
\end{itemize}

\textbf{Added}
\begin{itemize}
  \item §\emph{Two numbers, not four: the pair and the octet} --- the capacities built from the antipodal pair and
    the $4{+}4$ octet, with the closure argument, the $4{+}3$ step, the noble-gas statement, the
    direction on the doubling, and Linnett credited. Closure is stated as exhaustion of a row's
    capacity and not as a completed outer octet, which is false from the fourth row on: krypton's
    outermost eight is not $4{+}4$ under either filling rule, nor in the orbital account.
  \item §\emph{Aufbau: how the shell structure forms} --- the build-up as what makes energy minimisation select
    correctly, with the monotonicity measurement and the seed ladder (R2.2).
  \item A fourth-row table, $Z=19$--$36$ at $N=800$: mean $|$error$|$ $3.3\%$, all eighteen under-bound.
  \item Table 1 recomputed with no per-element constant: mean $|$error$|$ $1.6\%$, maximum $4.5\%$.
  \item §\emph{The shape of the computed field: constant charge per unit radius} --- the one scaling
    the construction does not contain: the charge per unit radius is the same in every shell and equal
    to $Z/2$, which fixes the atomic size near $2$~a.u.\ independently of $Z$ and so returns the
    near-constancy of atomic radii as a consequence rather than an input.
  \item A measurement bearing on the doubling: the geometric capacity $4\pi(r/d)^2$ per shell, on four
    atoms and three candidate fillings. The pair-and-octet filling is the more self-consistent, but on
    three informative shells only, so it is recorded as an anticipation of the filling at larger $Z$
    and not as a derivation of the doubling, which remains the principal open item.
  \item §\emph{Electronegativity from the shell energies} --- the Allen configuration energy, replacing ionization
    as the valence observable (R2.3).
  \item An analysis of what the spherical reduction costs, with the intra-shell repulsion table (R2.8).
  \item Three figures of the three-dimensional build-up, none in the original.
  \item The speculation caveat on ontological claims (R1.1), the molecular limits (R1.3), the basis-set
    concession (R1.4), and the imaginary-time clarification (R1.5).
\end{itemize}

\textbf{Restated}
\begin{itemize}
  \item Scope in its two senses: the \emph{computed} energies cover periods 1--4, hydrogen to krypton,
    all-electron and parameter-free; the \emph{account of the build-up} covers every row length and every
    closed shell. What is bounded at argon is the outward build-up, not the calculation.
  \item The shell structure as an output of the build-up, not of energy minimisation only.
  \item The Aufbau result as established for Li--C and open from nitrogen.
  \item The title, extended to name the new content: \emph{A 3D Multiphase Continuum Model for Atoms and the Periodic Table}. The model is three-dimensional; the periodic-table results are computed with the
    spherically homogenised reduction of it, as the text states throughout.
\end{itemize}

The revision retracts what it could not defend and we have tried to save none of it. What replaces it
is the construction: the table assembled from an antipodal pair and a dual-tetrahedral octet, every
row length and every closure following from what those two units can and cannot compose. That is a
different kind of claim from the original's --- checkable by composition rather than weighed against a
fitted number --- and we would ask that the manuscript be judged on it.

\end{document}
